Multi-touch technology allows one to control multiple simultaneous input points through the use of one's fingers. This provides a more natural form of interaction when compared to traditional input mediums, such as a keyboard and mouse. One area that may benefit greatly from this is computer gaming. Games, unlike most other software, are primarily used for their intrinsic enjoyment value. Thus, given the relatively unexplored nature of this area, valuable information may be gained through its exploration.

The effect of multi-touch interaction techniques on game enjoyment was investigated by comparing direct manipulation and gestures, that is the use of symbolic forms in order to relay scope and commands, as well as a combination of the two styles. In addition, the influence of game genre and previous gaming and touch-device experience were also analysed. For the overall project, a system consisting of three games of different genres, as well as a gesture recognition unit, was implemented. This document details the design and implementation of a puzzle game as well as the feature extraction unit necessary to reduce the dimensionality of data for further analysis by the gesture recognition engine.

User experimentation served to assess which interface produced the most flow. Flow, which describes the intrinsic enjoyment present in a task, was assessed using the Flow State Scale. The gestural interface was found to elicit more flow than either of the other two interfaces, while the mixed interface elicited more flow than the direct manipulation version. Genre and previous gaming and touch-device experience were not shown to have a significant effect on interface preference. However, previous gaming experience increased the overall level of flow experienced in all conditions.

The results seem to indicate that for games, the inclusion of gestural controls, where their use is appropriate, when compared to other interaction techniques, may increase the degree of flow, and thus overall game enjoyment, experienced by players.


